\section{Introduction}


The Ising model is a fundamental model in statistical mechanics used to study phase transitions and critical phenomena.
Despite its seemingly simple definition, including binary spin variables and nearest-neighbor interactions, it captures the essential features of spontaneous symmetry breaking and critical behaviour observed in even more complex systems.
In two dimensions, the Ising model exhibits a well-known second-order phase transition at a critical temperature $T_c$, where various physical quantities such as magnetization, susceptibility, and correlation length display power-law divergences characterized by critical exponents.

To study both static and dynamic properties of the Ising model, Monte Carlo simulations based on Markov-chain sampling are widely used.
Near the critical point, these algorithms are subject to \textbf{critical slowing down}, where the relaxation and autocorrelation time increase rapidly as the correlation length diverges.
In practice, this means that successive configurations generated by the simulation become highly correlated, leading to inefficient sampling and large statistical errors.
Understanding this is essential for accurate numerical studies close to $T_c$.

This study investigates the effective dynamics of the 2D Ising model as generated by different Monte Carlo update algorithms: the local Metropolis algorithm, which performs local single-spin updates, and the global Wolff cluster algorithm, which employs non-local cluster flips.
Although both algorithms sample from the same equilibrium distribution, they impose different dynamical behaviors on the system, leading to distinct scaling behaviours of the autocorrelation time near the critical point.
This scaling is characterized by the \textbf{dynamic critical exponent} $z$, which determines the divergence of time scales according to $\tau \sim \xi^{z}$ as the correlation length $\xi \to \infty$ at criticality.

The goals of this report are:

\begin{enumerate}
    \item To derive and discuss the theoretical background of dynamic critical behaviour and the definition of the dynamic exponent \(z\).
    \item To measure autocorrelation times of magnetization and energy for both Metropolis and Wolff dynamics as the temperature approaches $T_c$.
    \item To compare the scaling and efficiency of the two algorithms in the critical region.
\end{enumerate}

Section 2 introduces the theoretical framework, including the Ising model, Monte Carlo dynamics, and the concept of dynamic scaling. Section 3 outlines the simulation methodology and measurement procedures. Section 4 presents the numerical results and analysis, while Section 5 discusses the conclusions of the findings.